% Notes of the Montreal seminar, 1970 (cote 129, typed part): archive cote 129, site page 43 (PDF page 44).
% typescript. Draft compared: none.
% AI-transcribed by Claude Opus 5.5 (claude-opus-5-5, adaptive thinking, effort high), Message Batches API, batch mtl129_full (msgbatch_01AGnDKHwy6ATCM7ehUa6263), 2026-11-13; not yet checked by eye.
\scanpage{129/43}
\mspage{II.15}

L'action de $F$ et $V$ sur $D^*(G)$ se définit en posant,
pour tout $u \in D^*(G)$
\[ F\cdot u = F_{\underset{\rightarrow}{W}} \circ u \]
\[ V\cdot u = V_{\underset{\rightarrow}{W}} \circ u \]
où $F_{\underset{\rightarrow}{W}}$ et $V_{\underset{\rightarrow}{W}}$ sont les endomorphismes de Frobénius et de Verschiebung
\margin{XX} de $\underset{\rightarrow}{\mathbb{W}}$ , définis par passage à la limite des morphismes correspondant\hand{s} sur \inserted{les} $W_n$ .

On vérifie immédiatement, en revenant à la définition de l'ac-%
tion de $W$ sur $\underset{\rightarrow}{W}$ et en utilisant les propriétés de $F_{\underset{\rightarrow}{W}}$ et $V_{\underset{\rightarrow}{W}}$ , que
ces actions font de $D^*(G)$ un $\mathcal{D}$-module.

4.3.2. Montrons que l'on a un isomorphisme canonique
\[ D^*(G^{(p)}) \overset{\sim}{\rightarrow} (D^*(G))^{\sigma} \]
tel que
\[ D^*(\mathbb{F}_{G/k}) = F_{D^*(G)} \]
\[ D^*(\mathbb{V}_{G}) = V_{D^*(G)} \quad . \]

En effet, $k$ étant parfait, le changement de base $k \overset{\sigma}{\rightarrow} k$
induit un isomorphisme $\sigma$-linéaire
\[ \mathrm{Hom}(G, \mathbb{W}) \overset{\sim}{\rightarrow} \mathrm{Hom}(G^{(p)}, \mathbb{W}^{(p)}) \quad . \]

Mais $\underset{\rightarrow}{\mathbb{W}}$ étant défini sur le corps premier $\mathbb{F}_p$ , on a $\underset{\rightarrow}{\mathbb{W}}^{(p)} = \underset{\rightarrow}{\mathbb{W}}$
(car la formation de $\underset{\rightarrow}{\mathbb{W}}^{(p)}$ commute aux changements de base, cf. I, 5.2).
Par ailleurs
\[ D^*(\mathbb{F}_{G/k}) = \mathrm{Hom}(G^{(p)}, \underset{\rightarrow}{\mathbb{W}}) \rightarrow \mathrm{Hom}(G, \underset{\rightarrow}{\mathbb{W}}) \]
\margin{$\times$ \quad $\mapsto$}
\[ u \mapsto u \circ \mathbb{F}_{G/k} \]

\margin{43}
